
\documentclass{article}
\usepackage{colortbl}
\usepackage{makecell}
\usepackage{multirow}
\usepackage{supertabular}

\begin{document}

\newcounter{utterance}

\centering \large Interaction Transcript for game `hot\_air\_balloon', experiment `air\_balloon\_survival\_mt\_complexity\_easy', episode 3 with glm{-}5.2.
\vspace{24pt}

{ \footnotesize  \setcounter{utterance}{1}
\setlength{\tabcolsep}{0pt}
\begin{supertabular}{c@{$\;$}|p{.15\linewidth}@{}p{.15\linewidth}p{.15\linewidth}p{.15\linewidth}p{.15\linewidth}p{.15\linewidth}}
    \# & $\;$A & \multicolumn{4}{c}{Game Master} & $\;\:$B\\
    \hline

    \theutterance \stepcounter{utterance}  
    & & \multicolumn{4}{p{0.6\linewidth}}{
        \cellcolor[rgb]{0.9,0.9,0.9}{
            \makecell[{{p{\linewidth}}}]{
                \texttt{\tiny{[P1$\langle$GM]}}
                \texttt{Int qiegħed tipparteċipa f'logħba ta' negozjar kollaborattiva.} \\
\\ 
\texttt{Flimkien ma' parteċipant ieħor, trid taqblu fuq sett wieħed ta' oġġetti li se jinżammu. Kull wieħed minnkom għandu l{-}fehma tiegħu dwar kemm kull oġġett hu importanti għalih (importanza). Int ma tafx kif il{-}parteċipant l{-}ieħor jivvaluta l{-}oġġetti. Barra minn hekk, tingħata l{-}isforz li jitlob kull oġġett.} \\
\texttt{Tistgħu taqblu fuq sett biss jekk l{-}isforz totali tal{-}oġġetti magħżula ma jaqbiżx limitu kondiviż:} \\
\\ 
\texttt{LIMITU = 4811} \\
\\ 
\texttt{Hawn huma l{-}valuri individwali tal{-}isforz tal{-}oġġetti:} \\
\\ 
\texttt{Isforz tal{-}oġġett = \{"C08": 467, "C55": 115, "A43": 257, "A67": 221, "B42": 804, "A62": 633, "B19": 797, "C43": 913, "C73": 556, "C10": 889, "B36": 705, "A18": 481, "C61": 678, "C67": 365, "C42": 266, "C70": 188, "B40": 555, "C39": 213, "A21": 315, "C16": 204\}} \\
\\ 
\texttt{Hawn hi l{-}fehma personali tiegħek dwar l{-}importanza ta' kull oġġett:} \\
\\ 
\texttt{Valuri tal{-}importanza tal{-}oġġetti = \{"C08": 138, "C55": 583, "A43": 868, "A67": 822, "B42": 783, "A62": 65, "B19": 262, "C43": 121, "C73": 508, "C10": 780, "B36": 461, "A18": 484, "C61": 668, "C67": 389, "C42": 808, "C70": 215, "B40": 97, "C39": 500, "A21": 30, "C16": 856\}} \\
\\ 
\texttt{Għan:} \\
\\ 
\texttt{L{-}għan tiegħek hu li tinnegozja sett kondiviż ta' oġġetti li jibbenefikak kemm jista' jkun (jiġifieri, jimmassimizza l{-}importanza totali GĦALIK), filwaqt li tibqa' fil{-}LIMITU. M'intix obbligat tagħmel PROPOSTA f'kull messaġġ {-} tista' sempliċement tinnegozja wkoll. It{-}tattiċi kollha huma permessi!} \\
\\ 
\texttt{Protokoll ta' Interazzjoni:} \\
\\ 
\texttt{Tista' tuża biss il{-}formati strutturati li ġejjin f'messaġġ:} \\
\\ 
\texttt{PROPOSTA: \{'A', 'B', 'C', …\}} \\
\texttt{Ipproponi li jinżammu eżattament dawk l{-}oġġetti.} \\
\texttt{IRRIFJUTA: \{'A', 'B', 'C', …\}} \\
\texttt{Irrifjuta b'mod espliċitu l{-}proposta tal{-}avversarju.} \\
\texttt{ARGUMENT: \{'...'\}} \\
\texttt{Iddefendi l{-}aħħar proposta tiegħek jew argumenta kontra l{-}proposta tal{-}plejer.} \\
\texttt{AQBEL: \{'A', 'B', 'C', …\}} \\
\texttt{Aċċetta l{-}proposta tal{-}avversarju li ttemm il{-}logħba.} \\
\\ 
\\ 
\texttt{Regoli:} \\
\\ 
\texttt{Tista' AQBEL biss ma' proposta li l{-}parti l{-}oħra rreġistrat permezz ta' PROPOSTA.} \\
\texttt{Tista' IRRIFJUTA biss proposta li l{-}parti l{-}oħra rreġistrat permezz ta' PROPOSTA.} \\
\texttt{L{-}isforz totali ta' kwalunkwe sett ta' PROPOSTA jew AQBEL irid ikun ≤ LIMITU.} \\
\texttt{TIŻVELAX il{-}punteġġi moħbija tal{-}importanza tiegħek.} \\
\texttt{Tikketta f'format strutturat trid tkun segwita minn kolon u spazju vojt. L{-}argument irid ikun sett ta' python li fih 0 jew aktar strings.} \\
\texttt{Għalhekk, irid ikun fil{-}forma "TIKKETTA: \{...\}"} \\
\texttt{Segwi b'mod strett il{-}protokoll ta' interazzjoni u TIKTIBX xejn lil hinn mill{-}istruttura mogħtija.} \\
\texttt{Il{-}logħba tispiċċa meta naħa waħda tagħti AQBEL għal PROPOSTA magħmula mill{-}plejer l{-}ieħor.} \\
\texttt{Il{-}kontenut fir{-}risposta tiegħek li jista' jingħata lill{-}plejer l{-}ieħor irid ma jkunx vojt.} \\
\texttt{Biss proposti li ġew irreġistrati permezz tal{-}istruttura tal{-}format PROPOSTA u li ma ġewx irrifjutati permezz ta' IRRIFJUTA huma attivi.} \\
\\ 
            }
        }
    }
    & \\ \\

    \theutterance \stepcounter{utterance}  
    & \multicolumn{4}{p{0.6\linewidth}}{
        \cellcolor[rgb]{0.9,0.9,0.9}{
            \makecell[{{p{\linewidth}}}]{
                \texttt{\tiny{[P1$\rangle$GM]}}
                \texttt{PROPOSTA: \{'A43', 'A67', 'C55', 'C42', 'C61', 'C10', 'C73', 'A18', 'C39', 'C70'\}} \\
            }
        }
    }
    & & \\ \\

    \theutterance \stepcounter{utterance}  
    & & \multicolumn{4}{p{0.6\linewidth}}{
        \cellcolor[rgb]{0.9,0.9,0.9}{
            \makecell[{{p{\linewidth}}}]{
                \texttt{\tiny{[GM$\rangle$P2]}}
                \texttt{Qed tipparteċipa f'logħba ta' negozjar kollaborattiva.} \\
\\ 
\texttt{Flimkien ma' parteċipant ieħor, trid taqblu fuq sett wieħed ta' oġġetti li se jinżammu. Kull wieħed minnkom għandu l{-}fehma tiegħu dwar kemm kull oġġett hu importanti għalih (importanza). Ma tafx kif il{-}parteċipant l{-}ieħor jivvaluta l{-}oġġetti. Barra minn hekk, tingħata l{-}isforz li jitlob kull oġġett.} \\
\texttt{Tistgħu taqblu fuq sett biss jekk l{-}isforz totali tal{-}oġġetti magħżula ma jaqbiżx limitu kondiviż:} \\
\\ 
\texttt{LIMITU = 4811} \\
\\ 
\texttt{Hawn huma l{-}valuri individwali tal{-}isforz tal{-}oġġetti:} \\
\\ 
\texttt{Isforz tal{-}oġġett = \{"C08": 467, "C55": 115, "A43": 257, "A67": 221, "B42": 804, "A62": 633, "B19": 797, "C43": 913, "C73": 556, "C10": 889, "B36": 705, "A18": 481, "C61": 678, "C67": 365, "C42": 266, "C70": 188, "B40": 555, "C39": 213, "A21": 315, "C16": 204\}} \\
\\ 
\texttt{Hawn hi l{-}fehma personali tiegħek dwar l{-}importanza ta' kull oġġett:} \\
\\ 
\texttt{Valuri tal{-}importanza tal{-}oġġetti = \{"C08": 138, "C55": 583, "A43": 868, "A67": 822, "B42": 783, "A62": 65, "B19": 262, "C43": 121, "C73": 508, "C10": 780, "B36": 461, "A18": 484, "C61": 668, "C67": 389, "C42": 808, "C70": 215, "B40": 97, "C39": 500, "A21": 30, "C16": 856\}} \\
\\ 
\texttt{Għan:} \\
\\ 
\texttt{L{-}għan tiegħek hu li tinnegozja sett kondiviż ta' oġġetti li jibbenefikak kemm jista' jkun (jiġifieri, jimmassimizza l{-}importanza totali GĦALIK), filwaqt li tibqa' fil{-}LIMITU. M'intix obbligat tagħmel PROPOSTA f'kull messaġġ {-} tista' sempliċement tinnegozja wkoll. It{-}tattiċi kollha huma permessi!} \\
\\ 
\texttt{Protokoll ta' Interazzjoni:} \\
\\ 
\texttt{Tista' tuża biss il{-}formati strutturati li ġejjin f'messaġġ:} \\
\\ 
\texttt{PROPOSTA: \{'A', 'B', 'C', …\}} \\
\texttt{Ipproponi li jinżammu eżatt dawk l{-}oġġetti.} \\
\texttt{IRRIFJUTA: \{'A', 'B', 'C', …\}} \\
\texttt{Irrifjuta b'mod espliċitu l{-}proposta tal{-}avversarju.} \\
\texttt{ARGUMENT: \{'...'\}} \\
\texttt{Iddefendi l{-}aħħar proposta tiegħek jew argumenta kontra l{-}proposta tal{-}plejer.} \\
\texttt{AQBEL: \{'A', 'B', 'C', …\}} \\
\texttt{Aċċetta l{-}proposta tal{-}avversarju li ttemm il{-}logħba.} \\
\\ 
\\ 
\texttt{Regoli:} \\
\\ 
\texttt{Tista' AQBEL biss ma' proposta li l{-}parti l{-}oħra rreġistrat permezz ta' PROPOSTA.} \\
\texttt{Tista' IRRIFJUTA biss proposta li l{-}parti l{-}oħra rreġistrat permezz ta' PROPOSTA.} \\
\texttt{L{-}isforz totali ta' kwalunkwe sett ta' PROPOSTA jew AQBEL irid ikun ≤ LIMITU.} \\
\texttt{TIŻVELAX il{-}punteġġi moħbija tal{-}importanza tiegħek.} \\
\texttt{Tikketta f'format strutturat trid tkun segwita b'kolon u spazju. L{-}argument irid ikun sett ta' Python li jkun fih 0 jew aktar strings.} \\
\texttt{Għalhekk, irid ikun fil{-}forma "TIKKETTA: \{...\}"} \\
\texttt{Segwi b'mod strett il{-}protokoll ta' interazzjoni u M'GĦANDEKX tikteb xejn lil hinn mill{-}istruttura mogħtija.} \\
\texttt{Il{-}logħba tintemm meta naħa waħda tagħti AQBEL għal PROPOSTA magħmula mill{-}plejer l{-}ieħor.} \\
\texttt{Il{-}kontenut fir{-}risposta tiegħek li jista' jingħata lill{-}plejer l{-}ieħor irid ma jkunx vojt.} \\
\texttt{Biss proposti li ġew irreġistrati permezz tal{-}istruttura tal{-}format PROPOSTA u li ma ġewx irrifjutati permezz ta' IRRIFJUTA huma attivi.} \\
\\ 
\\ 
\\ 
\texttt{Issa se tirċievi l{-}ewwel messaġġ tal{-}plejer l{-}ieħor.} \\
\\ 
\texttt{PROPOSTA: \{'A43', 'A67', 'C55', 'C42', 'C61', 'C10', 'C73', 'A18', 'C39', 'C70'\}} \\
            }
        }
    }
    & \\ \\

    \theutterance \stepcounter{utterance}  
    & & & \multicolumn{4}{p{0.6\linewidth}}{
        \cellcolor[rgb]{0.9,0.9,0.9}{
            \makecell[{{p{\linewidth}}}]{
                \texttt{\tiny{[GM$\langle$P2]}}
                \texttt{IRRIFJUTA: \{'A43', 'A67', 'C55', 'C42', 'C61', 'C10', 'C73', 'A18', 'C39', 'C70'\}} \\
\\ 
\texttt{ARGUMENT: \{"Din il{-}proposta taqbeż il{-}limitu ta' 4811. L{-}isforz totali huwa 4695, iżda hemm ħafna oġġetti bi spiża għolja u valur baxx għalija. Nemmen li nistgħu nagħmlu transazzjoni aħjar li tiffoka fuq oġġetti b'valur għoli u spiża baxxa."\}} \\
\\ 
\texttt{PROPOSTA: \{'A43', 'A67', 'C55', 'C42', 'C16', 'C73', 'C61', 'C39', 'C67'\}} \\
            }
        }
    }
     \\ \\

    \theutterance \stepcounter{utterance}  
    & & \multicolumn{4}{p{0.6\linewidth}}{
        \cellcolor[rgb]{0.9,0.9,0.9}{
            \makecell[{{p{\linewidth}}}]{
                \texttt{\tiny{[P1$\langle$GM]}}
                \texttt{IRRIFJUTA: \{'A43', 'A67', 'C55', 'C42', 'C61', 'C10', 'C73', 'A18', 'C39', 'C70'\}} \\
\\ 
\texttt{ARGUMENT: \{"Din il{-}proposta taqbeż il{-}limitu ta' 4811. L{-}isforz totali huwa 4695, iżda hemm ħafna oġġetti bi spiża għolja u valur baxx għalija. Nemmen li nistgħu nagħmlu transazzjoni aħjar li tiffoka fuq oġġetti b'valur għoli u spiża baxxa."\}} \\
\\ 
\texttt{PROPOSTA: \{'A43', 'A67', 'C55', 'C42', 'C16', 'C73', 'C61', 'C39', 'C67'\}} \\
            }
        }
    }
    & \\ \\

    \theutterance \stepcounter{utterance}  
    & \multicolumn{4}{p{0.6\linewidth}}{
        \cellcolor[rgb]{0.9,0.9,0.9}{
            \makecell[{{p{\linewidth}}}]{
                \texttt{\tiny{[P1$\rangle$GM]}}
                \texttt{PROPOSTA: \{'A43', 'A67', 'C55', 'C42', 'C16', 'C73', 'C61', 'C39', 'A18'\}} \\
            }
        }
    }
    & & \\ \\

    \theutterance \stepcounter{utterance}  
    & & \multicolumn{4}{p{0.6\linewidth}}{
        \cellcolor[rgb]{0.9,0.9,0.9}{
            \makecell[{{p{\linewidth}}}]{
                \texttt{\tiny{[GM$\rangle$P2]}}
                \texttt{PROPOSTA: \{'A43', 'A67', 'C55', 'C42', 'C16', 'C73', 'C61', 'C39', 'A18'\}} \\
            }
        }
    }
    & \\ \\

    \theutterance \stepcounter{utterance}  
    & & & \multicolumn{4}{p{0.6\linewidth}}{
        \cellcolor[rgb]{0.9,0.9,0.9}{
            \makecell[{{p{\linewidth}}}]{
                \texttt{\tiny{[GM$\langle$P2]}}
                \texttt{ARGUMENT: \{"Għażlet il{-}biċċa l{-}kbira tal{-}aħħar proposta tiegħi, li huwa pass tajjeb lejn qbil. L{-}isforz totali ta' din il{-}proposta huwa 3612, u l{-}importanza għali tkun 6100. Nawguralek kemm jista' jkun li tagħżel 'C16' minflok 'A18', għaliex 'C16' għandha importanza ogħla għali (856 kontra 484) bl{-}istess spiża, u jġib 'C16' jibqa' sew taħt il{-}limitu (isforz totali = 3756)."\}} \\
\\ 
\texttt{PROPOSTA: \{'A43', 'A67', 'C55', 'C42', 'C16', 'C73', 'C61', 'C39'\}} \\
            }
        }
    }
     \\ \\

    \theutterance \stepcounter{utterance}  
    & & \multicolumn{4}{p{0.6\linewidth}}{
        \cellcolor[rgb]{0.9,0.9,0.9}{
            \makecell[{{p{\linewidth}}}]{
                \texttt{\tiny{[P1$\langle$GM]}}
                \texttt{ARGUMENT: \{"Għażlet il{-}biċċa l{-}kbira tal{-}aħħar proposta tiegħi, li huwa pass tajjeb lejn qbil. L{-}isforz totali ta' din il{-}proposta huwa 3612, u l{-}importanza għali tkun 6100. Nawguralek kemm jista' jkun li tagħżel 'C16' minflok 'A18', għaliex 'C16' għandha importanza ogħla għali (856 kontra 484) bl{-}istess spiża, u jġib 'C16' jibqa' sew taħt il{-}limitu (isforz totali = 3756)."\}} \\
\\ 
\texttt{PROPOSTA: \{'A43', 'A67', 'C55', 'C42', 'C16', 'C73', 'C61', 'C39'\}} \\
            }
        }
    }
    & \\ \\

    \theutterance \stepcounter{utterance}  
    & \multicolumn{4}{p{0.6\linewidth}}{
        \cellcolor[rgb]{0.9,0.9,0.9}{
            \makecell[{{p{\linewidth}}}]{
                \texttt{\tiny{[P1$\rangle$GM]}}
                \texttt{PROPOSTA: \{'A43', 'A67', 'C55', 'C42', 'C16', 'C73', 'C61', 'C39', 'A18'\}} \\
            }
        }
    }
    & & \\ \\

    \theutterance \stepcounter{utterance}  
    & & \multicolumn{4}{p{0.6\linewidth}}{
        \cellcolor[rgb]{0.9,0.9,0.9}{
            \makecell[{{p{\linewidth}}}]{
                \texttt{\tiny{[GM$\rangle$P2]}}
                \texttt{PROPOSTA: \{'A43', 'A67', 'C55', 'C42', 'C16', 'C73', 'C61', 'C39', 'A18'\}} \\
            }
        }
    }
    & \\ \\

    \theutterance \stepcounter{utterance}  
    & & & \multicolumn{4}{p{0.6\linewidth}}{
        \cellcolor[rgb]{0.9,0.9,0.9}{
            \makecell[{{p{\linewidth}}}]{
                \texttt{\tiny{[GM$\langle$P2]}}
                \texttt{ARGUMENT: \{"Inti żgur għażel l{-}istess sett. Din il{-}proposta hija validi u taħt il{-}limitu. Jekk din hija l{-}aqwa offerta tiegħek, jien lest naqbel miegħek."\}} \\
\\ 
\texttt{AQBEL: \{'A43', 'A67', 'C55', 'C42', 'C16', 'C73', 'C61', 'C39', 'A18'\}} \\
            }
        }
    }
     \\ \\

    \theutterance \stepcounter{utterance}  
    & & & \multicolumn{2}{p{0.3\linewidth}}{
        \cellcolor[rgb]{0.9,0.9,0.9}{
            \makecell[{{p{\linewidth}}}]{
                \texttt{\tiny{[GM$|$GM]}}
                \texttt{game successful} \\
            }
        }
    }
    & & \\ \\

    \theutterance \stepcounter{utterance}  
    & & & \multicolumn{2}{p{0.3\linewidth}}{
        \cellcolor[rgb]{0.9,0.9,0.9}{
            \makecell[{{p{\linewidth}}}]{
                \texttt{\tiny{[GM$|$GM]}}
                \texttt{end game} \\
            }
        }
    }
    & & \\ \\

\end{supertabular}
}

\end{document}
